% TOP_PROBLEM: {"id": 30000465, "title": "Continuous-Spectrum Analysis of Nelson-Type Models", "category_id": 16, "category": {"id": 16, "name": "physics", "display_name": "Mathematical Physics", "description": "Problems at the intersection of mathematics and physics.", "slug": "physics", "order_index": 16, "created_at": "2026-07-31T15:26:25.671Z"}, "status": "partially_solved", "statusQualification": "The curated QED energy set reverses the inequality in primary Remark 6. The report specifies (sigma(H) minus sigma_pp(H)) intersected with energies below Sigma_el; the catalogue instead removes the interval below Sigma_el. Both are displayed explicitly and cannot be treated as equivalent."}
% FIRST_POSTED: 2026-10-11
% ENTRY_KIND: statement_only
% SUBMISSION: {"kind": "statement_only", "category": "top_problems", "problem_id": "30000465", "proposal_type": "established", "source": "The Rigorous Renormalization Group (2006), Oberwolfach Reports, DOI 10.4171/owr/2006/17, pp. 1027--1076; official report source and mathematical context. https://doi.org/10.4171/owr/2006/17; Official EMS Press PDF of the same Oberwolfach report; PDF page 40, printed page 1066, Remark 6. https://ems.press/content/serial-article-files/46049; Møller--Schenker--Vugalter, massive Nelson spectrum; reference in the saved dated catalogue literature review. https://arxiv.org/abs/1207.6729", "status_evidence": "Saved source status is partially_solved; retained unchanged. The official historical OWR source was inspected on 2026-10-11. The saved literature assessment is dated 2026-08-21 and is not a fresh comprehensive status verification. ", "status_checked_on": "2026-10-11", "duplicate_check": "Checked canonical identity and normalized formulations against the complete saved UnsolvedMath registry; No exact normalized duplicate found. Mathematical equivalence and current literature were not exhaustively searched; no new ID allocated.", "attribution": "Israel Michael Sigal, reporting joint work with Jürg Fröhlich and Marcel Griesemer. Statement preparation and source audit assisted by OpenAI GPT-6 (Codex); no human mathematical authorship or independent proof verification claimed."}
% SOURCE_STATEMENT: "Can the limiting-absorption and local-decay method be extended to $\\sigma(H)\\setminus\\sigma_{\\mathrm{pp}}(H)$ for the Nelson model and to $\\sigma(H)\\setminus(\\sigma_{\\mathrm{pp}}(H)\\cup(-\\infty,\\Sigma_{\\mathrm{el}}))$ for the standard model of nonrelativistic QED, where $\\Sigma_{\\mathrm{el}}=\\inf\\sigma(H_{\\mathrm{el}})$?"
% Preparation and final self-audit: CONTRIBUTING.md prompts C and D; source/audit artifacts retained in the local physics_statements_2026_10_11 batch.
% Imported from: physics_statements.tex; UnsolvedMath ID 30000465
% !TeX program = lualatex
\documentclass[11pt]{article}
\usepackage[margin=25mm]{geometry}
\usepackage{fontspec}
\setmainfont{Latin Modern Roman}
\usepackage{amsmath,amssymb,mathtools}
\usepackage{unicode-math}
\setmathfont{Latin Modern Math}
\usepackage{xurl,hyperref}
\hypersetup{colorlinks=true,urlcolor=blue}
\setcounter{secnumdepth}{0}
\setlength{\parindent}{0pt}
\setlength{\parskip}{5pt}
\setlength{\emergencystretch}{3em}
\newcommand{\sref}[2]{\hyperlink{source:#1:#2}{[S#2]}}
\newcommand{\cataloguescope}[1]{\paragraph{Recorded target.}#1}
\newcommand{\R}{\mathbb R}
\newcommand{\C}{\mathbb C}
\newcommand{\N}{\mathbb N}
\newcommand{\Z}{\mathbb Z}
\newcommand{\Q}{\mathbb Q}
\newcommand{\F}{\mathbb F}
\newcommand{\eps}{\varepsilon}
\begin{document}
% UnsolvedMath ID 30000465; source catalogue code OWR-1266-004
\setcounter{section}{30000464}
\section{Continuous-Spectrum Analysis of Nelson-Type Models}\label{30000465}
{\small UnsolvedMath ID 30000465 · Mathematical Physics}
\subsection{Definitions and mathematical statement}

\paragraph{Shared notation.}$\mathbb N=\{1,2,\ldots\}$; the standard number fields have their usual meanings. All additional conventions are stated locally.


\paragraph{Definitions and conventions.} $H$ is the selfadjoint interacting Hamiltonian of the report's Nelson or nonrelativistic-QED model; $\sigma(H)$ is its spectrum and $\sigma_{\mathrm{pp}}(H)$ the set of its eigenvalues. $H_{\mathrm{el}}$ is the electronic Hamiltonian and $\Sigma_{\mathrm{el}}=\inf\sigma(H_{\mathrm{el}})$ is the threshold named in the report. A limiting-absorption principle establishes suitably weighted boundary values of $(H-E\mp i\varepsilon)^{-1}$ as $\varepsilon\downarrow0$ on continuous spectral energies; local-decay estimates control the corresponding time evolution.

\paragraph{Source scope and qualifications.} The curated QED energy set reverses the inequality in primary Remark 6. The report specifies (sigma(H) minus sigma\_pp(H)) intersected with energies below Sigma\_el; the catalogue instead removes the interval below Sigma\_el. Both are displayed explicitly and cannot be treated as equivalent.

\paragraph{Primary-source energy range.} Remark 6 requests $\sigma(H)\setminus\sigma_{\mathrm{pp}}(H)$ for the Nelson model and $\bigl(\sigma(H)\setminus\sigma_{\mathrm{pp}}(H)\bigr)\cap(-\infty,\Sigma_{\mathrm{el}})$ for the standard nonrelativistic-QED model. This is the source-faithful range, distinct from the preserved catalogue formula below.

\paragraph{Catalogue mathematical target.}
\paragraph{Catalogue wording retained for traceability.}
Can the limiting-absorption and local-decay method be extended to $\sigma(H)\setminus\sigma_{\mathrm{pp}}(H)$ for the Nelson model and to $\sigma(H)\setminus(\sigma_{\mathrm{pp}}(H)\cup(-\infty,\Sigma_{\mathrm{el}}))$ for the standard model of nonrelativistic QED, where $\Sigma_{\mathrm{el}}=\inf\sigma(H_{\mathrm{el}})$?

\paragraph{Saved literature context.} The supplied catalogue's assessment dated 2026-08-21 records: Later spectral work proves limiting-absorption and continuous-spectrum results for important Nelson-model regimes, including below multi-boson thresholds, but the full extension to the general energy intervals stated in the conjecture is not established by the targeted search. The assessment is retained as dated provenance; it is not a new comprehensive literature or open-status certification.
\paragraph{Scope and source qualification.}The curated QED energy set reverses the inequality in primary Remark 6. The report specifies (sigma(H) minus sigma_pp(H)) intersected with energies below Sigma_el; the catalogue instead removes the interval below Sigma_el. Both are displayed explicitly and cannot be treated as equivalent.
\subsection{Short English statement}
Extend the report's continuous-spectrum resolvent and decay method to its larger proposed energy ranges.
\subsection{Sources}
\begin{itemize}
\item[\hypertarget{source:30000465:1}{S1}] ULAM.AI and the UnsolvedMath Contributors, UnsolvedMath, record 30000465 (OWR-1266-004). Catalogue curation under CC BY 4.0; definitions and English exposition added locally.
\url{https://huggingface.co/datasets/ulamai/UnsolvedMath}.
\item[\hypertarget{source:30000465:2}{S2}] The Rigorous Renormalization Group (2006), Oberwolfach Reports, DOI 10.4171/owr/2006/17, pp. 1027--1076; official report source and mathematical context.
\url{https://doi.org/10.4171/owr/2006/17}.
\item[\hypertarget{source:30000465:3}{S3}] Official EMS Press PDF of the same Oberwolfach report; PDF page 40, printed page 1066, Remark 6.
\url{https://ems.press/content/serial-article-files/46049}.
\item[\hypertarget{source:30000465:4}{S4}] Møller--Schenker--Vugalter, massive Nelson spectrum; reference in the saved dated catalogue literature review.
\url{https://arxiv.org/abs/1207.6729}.
\end{itemize}
\label{end:30000465}
\end{document}
